% GENERATED FILE. Edit canonical lessons, then run npm run generate:books.
% canonical-source-hash: fnv1a64:e8cd4adc523d6521
% canonical-lessons: GE-C26-sitzen, GE-C26-stehen, GE-C26-schlafen, GE-C26-hoeren

\chapter{Sitting, Standing, Sleeping, Hearing}
\label{ch:ge-still-verbs}
\hlchaptermodality{26}

\begin{chapteropening}
\textbf{By the end of this chapter:} \emph{I can say I sit, I stand, I sleep and I hear in German, tell the second shift's t-branch from its p-branch, break a verb's a to ä for du and er, and say which of the four changes German made alone and which it made alongside English.}

The chapter closes on hören, the one verb of the four whose vowel cannot break because it was born umlauted: the reader produces all four ich forms and all four er forms, names schlafen as the only breaker, tells Grimm's law on the h of hören and Hund from the second shift on sitzen and schlafen, and hangs gern on the new verb. It reaches back past its own chapter to Chapter 17's Grimm's law, Chapter 22's Hund and Katze with the Hund etymology that had never been revisited, and Chapter 25's mögen, lieben and gern, all three of which were orphans until here.
\end{chapteropening}

\section[sitzen]{sitzen — the same word as \emph{sit}, respelt by one law}

\label{lesson:GE-C26-sitzen}



\begin{quote}
\textbf{Helfen} \emph{is} \textbf{help}, split by the second sound shift: Germanic \textbf{p} became German \textbf{f}. That law has other branches. This verb rides one.
\end{quote}

\begin{sounds}
\begin{itemize}
\raggedright
  \item \texttt{tz-affricate} — \textbf{tz} is one squeezed sound, \emph{ts}, as in English \emph{cats}.
  \item \texttt{i-short} — a short, tight \emph{i}: \textbf{sitzen} = \emph{ZIT-sen}, buzzed at the front because German \textbf{s} before a vowel is voiced.
\end{itemize}
\end{sounds}

\subsection*{You'll want to know: sitzen}
\textbf{sitzen} means \textquotedblleft{}\textbf{to sit}.\textquotedblright{}

\begin{itemize}
\raggedright
  \item \textbf{ich sitze} — I sit · \textbf{wir sitzen} — we sit
  \item \textbf{du sitzt} — you sit · \textbf{ihr sitzt} — you all sit
  \item \textbf{er sitzt} — he sits · \textbf{sie sitzen} — they sit
\end{itemize}

\begin{quote}
\textbf{Ich sitze.} — a whole sentence, two words.
\end{quote}

\begin{grammarlens}[title={Grammar Lens: strong does not mean it breaks}]
\emph{Sitzen} is a \textbf{strong} verb — its past is \emph{saß}, not \emph{sitzte}. But its present never breaks its vowel: \emph{du sitzt, er sitzt}. Beside it, \emph{helfen} breaks to \emph{du hilfst}.

So \textquotedblleft{}strong\textquotedblright{} is a fact about a verb's \textbf{past}; the \emph{du}/\emph{er} break is a separate habit only some strong verbs have.

Chapter 25 gave two ways to say you like something: \textbf{ich mag} with a thing, \textbf{gern} hung on a verb. This verb takes the second.

\begin{quote}
\textbf{Ich sitze gern.} — \textquotedblleft{}I like sitting.\textquotedblright{}
\end{quote}
\end{grammarlens}

\begin{cousinweb}
\textbf{sitzen} is Proto-Germanic \emph{\textbf{*sitjaną}}, and English \textbf{sit} is the same inherited word — Old English \emph{sittan}, Old High German \emph{sizzen}.

The second, High German shift moved Germanic \textbf{p} to \textbf{f} in \emph{helfen}. It moved Germanic \textbf{t} just as hard, to \textbf{z} (said \emph{ts}) or \textbf{ss}: \emph{sittan} $\to$ \textbf{sitzen}, \emph{water} $\to$ \textbf{Wasser}, \emph{eat} $\to$ \textbf{essen}, \emph{ten} $\to$ \textbf{zehn}, \emph{bite} $\to$ \textbf{beißen}. Chapter 11 handed you \emph{Wasser} and Chapter 6 handed you \emph{zehn} long before there was a name for what had happened to them.

Under both lies PIE \emph{\textbf{*sed-}}, \textquotedblleft{}to sit,\textquotedblright{} one of the family's great providers. Latin \emph{sedēre} gives \textbf{session}, \textbf{sedentary} and \textbf{preside}; Latin \emph{sedēs}, \textquotedblleft{}a seat,\textquotedblright{} gives the bishop's \textbf{see}. Greek \emph{hédra}, also \textquotedblleft{}a seat,\textquotedblright{} gives \textbf{cathedral} — the church with the bishop's chair — and \textbf{polyhedron}.
\end{cousinweb}

\subsection*{Your turn}
Say these aloud:

\begin{itemize}
\raggedright
  \item \textquotedblleft{}ich sitze, du sitzt, er sitzt\textquotedblright{} — the \emph{i} never moves
  \item the contrast — \textquotedblleft{}du sitzt\textquotedblright{} beside \textquotedblleft{}du hilfst\textquotedblright{}
  \item \textquotedblleft{}Ich sitze gern.\textquotedblright{}
  \item the \emph{t} branch four times — \textquotedblleft{}sitzen/sit, Wasser/water, essen/eat, zehn/ten\textquotedblright{}
  \item the \emph{p} branch beside it — \textquotedblleft{}helfen/help\textquotedblright{}
  \item the seat words — \textquotedblleft{}sitzen, sedēre, hédra, cathedral\textquotedblright{}
\end{itemize}

\begin{tcolorbox}[breakable,colback=teal!4,colframe=teal!35!black,title={Before you move on}]
Give the \emph{du} and \emph{er} forms of \emph{sitzen}. (\emph{\textbf{Du sitzt, er sitzt}} — no break.) Which verb from Chapter 25 does break, and how? (\emph{\textbf{Helfen}} — \emph{du hilfst}.) Which sound did the second shift turn German's \textbf{z} and \textbf{ss} out of? (An old \textbf{t}.) Name two pairs from that branch. (\textbf{Wasser/water}, \textbf{essen/eat}, \textbf{zehn/ten}.) Say you like sitting. (\emph{\textbf{Ich sitze gern}}.)
\end{tcolorbox}
\section[stehen]{stehen — the verb that was hiding inside \emph{verstehen}}

\label{lesson:GE-C26-stehen}



\begin{quote}
Chapter 24 broke \textbf{verstehen} into \emph{ver-} plus \emph{stehen}, and named \emph{stehen} as English \textbf{stand}. Here is that verb alone, and the letter the two languages disagree about.
\end{quote}

\begin{sounds}
\begin{itemize}
\raggedright
  \item \texttt{st-initial-scht} — \textbf{st} at the start of a German word is said \emph{sht}.
  \item \texttt{h-silent-lengthening} — the \textbf{h} is silent and marks the long vowel: \textbf{stehen} = \emph{SHTAY-en}.
\end{itemize}
\end{sounds}

\subsection*{You'll want to know: stehen}
\textbf{stehen} means \textquotedblleft{}\textbf{to stand}.\textquotedblright{}

\begin{itemize}
\raggedright
  \item \textbf{ich stehe} — I stand · \textbf{wir stehen} — we stand
  \item \textbf{du stehst} — you stand · \textbf{ihr steht} — you all stand
  \item \textbf{er steht} — he stands · \textbf{sie stehen} — they stand
\end{itemize}

Like \emph{sitzen}, it keeps its vowel all the way through. Two postures, two sentences:

\begin{quote}
\textbf{Ich sitze. Ich stehe.}
\end{quote}

\begin{cousinweb}
Germanic ran two stems off one root here: a short \emph{\textbf{*stāną}} and a longer \emph{\textbf{*standaną}} with an \textbf{n} wedged inside. English generalised the long one and got \textbf{stand}; German took the short one and got \textbf{stehen}, Old High German \emph{stān}.

German did not lose the \emph{n}. It kept it in the past: \textbf{ich stand}, and the participle \textbf{gestanden}. German's past tense is spelt the way English's present is, out of the same drawer of stems.

The root is PIE \emph{\textbf{*steh\textsubscript{2}-}}, the one Chapter 24 named, and it is everywhere. Latin \emph{stāre} gave \textbf{state}, \textbf{station}, \textbf{stable}, \textbf{constant}. Greek \emph{hístēmi}, \textquotedblleft{}I set up,\textquotedblright{} gave \textbf{static}, \textbf{system} and \textbf{ecstasy} — standing outside yourself. German built \textbf{die Stadt}, a city: a standing-place. English built \textbf{stead}, \textbf{steady} and \textbf{instead}.
\end{cousinweb}

\begin{culture}
Notice the shape of what has happened. \emph{Sitzen} goes back to \emph{\textbf{*sed-}}, and so does Latin's \emph{sedēre}. \emph{Stehen} goes back to \emph{\textbf{*steh\textsubscript{2}-}}, and so does Latin's \emph{stāre}.

Nothing was borrowed in either direction. Two families inherited the same two roots and held them to the same two meanings for thousands of years, because sitting and standing are things everybody does.
\end{culture}

\subsection*{Your turn}
Say these aloud:

\begin{itemize}
\raggedright
  \item \textquotedblleft{}ich stehe, du stehst, er steht\textquotedblright{} — \emph{SHTAY}, not \emph{STAY}
  \item both postures — \textquotedblleft{}Ich sitze. Ich stehe.\textquotedblright{}
  \item the verb inside the verb — \textquotedblleft{}stehen, verstehen\textquotedblright{}
  \item where the \emph{n} went — \textquotedblleft{}ich stehe, ich stand, gestanden\textquotedblright{}
  \item the Latin pair — \textquotedblleft{}sitzen/sedēre, stehen/stāre\textquotedblright{}
  \item what stands in English — \textquotedblleft{}state, station, stable, stead\textquotedblright{}
\end{itemize}

\begin{tcolorbox}[breakable,colback=teal!4,colframe=teal!35!black,title={Before you move on}]
Give the \emph{du} and \emph{er} forms. (\emph{\textbf{Du stehst, er steht}}.) Which letter does English's \emph{stand} carry that German's \emph{stehen} does not? (An \textbf{n}.) Where does German still keep it? (In the \textbf{past} — \emph{ich stand, gestanden}.) Which longer verb from Chapter 24 has \emph{stehen} inside it? (\emph{\textbf{Verstehen}}.) Say \textquotedblleft{}I sit. I stand.\textquotedblright{} (\emph{\textbf{Ich sitze. Ich stehe.}})
\end{tcolorbox}
\section[schlafen]{schlafen — \emph{sleep}, under two coats of paint}

\label{lesson:GE-C26-schlafen}



\begin{quote}
\emph{Sitzen} and \emph{stehen} held their vowels still. This one does not — and it does not break the way the Chapter 24 verbs broke either.
\end{quote}

\begin{sounds}
\begin{itemize}
\raggedright
  \item \texttt{sch-sh} — \textbf{sch} is English \emph{sh}: \textbf{schlafen} = \emph{SHLAH-fen}.
  \item \texttt{vowel-ae} — \textbf{ä} is a flat \emph{eh}: \textbf{schläft} = \emph{SHLEFT}.
\end{itemize}
\end{sounds}

\subsection*{You'll want to know: schlafen}
\textbf{schlafen} means \textquotedblleft{}\textbf{to sleep}.\textquotedblright{}

\begin{itemize}
\raggedright
  \item \textbf{ich schlafe} · \textbf{wir schlafen}
  \item \textbf{du schläfst} · \textbf{ihr schlaft}
  \item \textbf{er schläft} · \textbf{sie schlafen}
\end{itemize}

\begin{quote}
\textbf{Die Katze schläft.} — \textquotedblleft{}The cat is sleeping.\textquotedblright{}
\end{quote}

\begin{grammarlens}[title={Grammar Lens: the second kind of break}]
Chapter 24 gave you one break: long \emph{e} goes to \textbf{i} for \emph{du} and \emph{er} — \emph{du liest}, \emph{er hilft}. Here is the other: \textbf{a} takes two dots and becomes \textbf{ä} — \emph{du schläfst}, \emph{er schläft}.

Those two dots are the umlaut, the mark that turns \emph{Buch} into \emph{Bücher}. The break lands in the same slot as before, \textbf{du} and \textbf{er} and nobody else, so only its raw material is new: \emph{e} verbs brighten to \emph{i}, \emph{a} verbs to \emph{ä}.

\begin{quote}
\textbf{Es regnet. Die Katze schläft.}
\end{quote}
\end{grammarlens}

\begin{cousinweb}
\textbf{schlafen} is Proto-Germanic \emph{\textbf{*slēpaną}}, and English \textbf{sleep} is the same inherited word — Old English \emph{slæpan}, Old High German \emph{slāfan}.

Two German changes sit between them. The first you know: the second shift, Germanic \textbf{p} to German \textbf{f}, as in \emph{helfen}. The second is new — German turned \textbf{s} into \emph{sh} before \emph{l}, \emph{m}, \emph{n} and \emph{w}, and wrote it down: \emph{sleep} $\to$ \textbf{schlafen}, \emph{swim} $\to$ \textbf{schwimmen}, \emph{snow} $\to$ \textbf{Schnee}, \emph{smith} $\to$ \textbf{Schmied}. It did the same before \emph{p} and \emph{t} and never spelt it, which is why \emph{stehen} is said \emph{SHTAY-en}.

Now the honest part. The old Indo-European verb for sleeping was \emph{\textbf{*swep-}}, still in plain view: Latin \emph{somnus} gives \textbf{insomnia} and \textbf{somnolent}, Latin \emph{sopor} gives \textbf{soporific}, Greek \emph{hýpnos} gives \textbf{hypnosis}, and Old Norse \emph{sofa} is Swedish \emph{sova} today. German and English both threw that verb away and built a new one on \emph{\textbf{*sleb-}}, \textquotedblleft{}to be slack\textquotedblright{} — German still has \textbf{schlaff}, \textquotedblleft{}limp.\textquotedblright{} They made that swap \textbf{together}, before they were two languages.
\end{cousinweb}

\subsection*{Your turn}
Say these aloud:

\begin{itemize}
\raggedright
  \item \textquotedblleft{}ich schlafe, du schläfst, er schläft\textquotedblright{} — the \emph{a} brightens to \emph{ä}
  \item the three verbs — \textquotedblleft{}Ich sitze. Ich stehe. Ich schlafe.\textquotedblright{}
  \item \textquotedblleft{}Es regnet. Die Katze schläft.\textquotedblright{}
  \item both breaks in the same slot — \textquotedblleft{}du hilfst\textquotedblright{} beside \textquotedblleft{}du schläfst\textquotedblright{}
  \item the \emph{sch} row — \textquotedblleft{}schlafen, schwimmen, Schnee, Schmied\textquotedblright{}
  \item the verb they both threw away — \textquotedblleft{}somnus, hypnos, sova\textquotedblright{}
\end{itemize}

\begin{tcolorbox}[breakable,colback=teal!4,colframe=teal!35!black,title={Before you move on}]
Give the \emph{du} and \emph{er} forms. (\emph{\textbf{Du schläfst, er schläft}} — two dots on the \emph{a}.) Which two people does a German verb break for? (\textbf{Du} and \textbf{er}, always.) Name the two changes between \emph{sleep} and \emph{schlafen}. (\textbf{p} to \textbf{f}, and \textbf{s} to \textbf{sch}.) Why is \emph{stehen} said \emph{SHTAY-en}? (\textbf{Same change}, never spelt.) Say \textquotedblleft{}the cat is sleeping.\textquotedblright{} (\emph{\textbf{Die Katze schläft}}.)
\end{tcolorbox}
\section[hören]{hören — \emph{hear}, and the end of a chapter that never moved}

\label{lesson:GE-C26-hoeren}



\begin{quote}
\textbf{Ich sitze. Ich stehe. Ich schlafe.} Three verbs that stay put. Here is the fourth — the sense that works with your eyes shut.
\end{quote}

\begin{sounds}
\begin{itemize}
\raggedright
  \item \texttt{umlaut-oe} — \textbf{ö} is the \emph{e} of \emph{bed} said with rounded lips.
  \item \texttt{r-uvular} — the \textbf{r} is a soft scrape at the back: \textbf{hören} = \emph{HUR-en}.
\end{itemize}
\end{sounds}

\subsection*{You'll want to know: hören}
\textbf{hören} means \textquotedblleft{}\textbf{to hear}.\textquotedblright{}

\begin{itemize}
\raggedright
  \item \textbf{ich höre} · \textbf{wir hören}
  \item \textbf{du hörst} · \textbf{ihr hört}
  \item \textbf{er hört} · \textbf{sie hören}
\end{itemize}

The \textbf{ö} never moves. \emph{Schlafen} grew its dots for \emph{du} and \emph{er}; \emph{hören} was born with them, so nothing is left to break.

\begin{quote}
\textbf{Der Hund hört. Die Katze hört.}
\end{quote}

\begin{grammarlens}[title={Grammar Lens: liking it, the Chapter 25 way}]
\textbf{gern} hangs on this verb like any other:

\begin{quote}
\textbf{Ich höre gern.} — \textquotedblleft{}I like listening.\textquotedblright{}
\end{quote}

\emph{Gern} is English \textbf{yearn}: underneath, \textquotedblleft{}I hear yearningly.\textquotedblright{} Beside it is the other way of liking, \textbf{ich mag} — English \textbf{may}, with the power drained out.
\end{grammarlens}

\begin{cousinweb}
\textbf{hören} is Proto-Germanic \emph{\textbf{*hauzijaną}}, and English \textbf{hear} is the same inherited word.

This time the change is one they \textbf{share}. Gothic still says \emph{hausjan}, with an \textbf{s}; Old English \emph{hīeran} and Old High German \emph{hōren} had both turned it to \textbf{r}. \emph{Sitzen} shows where German left English; \emph{hören} shows a road they walked together.

The \textbf{h} is Grimm's law, the first shift: an old \textbf{k} became Germanic \textbf{h}. Greek kept the \emph{k} — \emph{akoúein}, \textquotedblleft{}to hear,\textquotedblright{} which English borrowed back as \textbf{acoustic}. Same swap that made \emph{*kwon-} into \textbf{Hund} while Latin kept \emph{canis}.

Some break the root up further, into \textquotedblleft{}sharp\textquotedblright{} plus \textquotedblleft{}ear\textquotedblright{} — which would make \emph{hören} \textquotedblleft{}sharp-eared\textquotedblright{} and put \textbf{das Ohr} inside it. Cited often, not settled.
\end{cousinweb}

\begin{culture}
\textbf{gehören} means \textquotedblleft{}\textbf{to belong}.\textquotedblright{} It is \emph{ge-} welded onto \emph{hören}, and began as \textquotedblleft{}to hearken to\textquotedblright{}: what belongs to you answers when called.
\end{culture}

\subsection*{Your turn}
Say these aloud:

\begin{itemize}
\raggedright
  \item all four — \textquotedblleft{}Ich sitze. Ich stehe. Ich schlafe. Ich höre.\textquotedblright{}
  \item the four \emph{er} forms — \textquotedblleft{}er sitzt, er steht, er schläft, er hört\textquotedblright{}
  \item \textquotedblleft{}Der Hund hört. Die Katze schläft.\textquotedblright{}
  \item \textquotedblleft{}Ich höre gern.\textquotedblright{}
  \item the shared change — \textquotedblleft{}hausjan, hīeran, hōren\textquotedblright{}
  \item the \emph{h} row — \textquotedblleft{}hören/akoúein, Hund/canis\textquotedblright{}
  \item \textquotedblleft{}gehören is \emph{belong}\textquotedblright{}
\end{itemize}

\begin{tcolorbox}[breakable,colback=teal!4,colframe=teal!35!black,title={Before you move on}]
Give all four \emph{ich} forms of this chapter. (\emph{\textbf{Ich sitze, ich stehe, ich schlafe, ich höre}}.) Which breaks its vowel? (\emph{\textbf{Schlafen}} — \emph{a} to \emph{ä}.) Which letter did Gothic keep where German and English have \emph{r}? (An \textbf{s}.) Which law put the \emph{h} on \emph{hören} and \emph{Hund}? (\textbf{Grimm's law}.)
\end{tcolorbox}
